\chapter{实现}

\section{RISC-V 与 AArch64 后端}

当前系统已完成 RISC-V 与 AArch64 两个主后端实现。RISC-V 路径支持压缩指令、ISA 字符串驱动的 \capstone\ mode 配置，以及面向扩展相关假阳性的 \code{--filter-ext} 逻辑。AArch64 路径则采用更严格的 mismatch 口径，用以减少“反汇编器可识别但用户态不可执行”的噪声样本。

RISC-V 的一个关键实现点是 raw hidden 与 filtered hidden 的双统计口径。当前代码在启用 \code{--filter-ext} 时同时保留 \code{raw\_hidden\_count}、过滤后 \code{hidden\_count}、以及按扩展分组的 \code{filtered\_hidden\_ext\_counts}。这样，同一轮 Layer A 实验既能保留原始异常候选，又能给出更接近当前 ISA 字符串视图的 triage 口径，避免为了比较 raw H 与 filtered H 而重复执行同一批编码。

AArch64 的实现选择则体现为另一种工程取舍：当前系统更强调严格 mismatch，而不是宽松 anomaly tally。原因在于，若把大量用户态不可执行但可被反汇编器识别的编码一律计入“异常”，评估章节会被噪声主导。与其追求更大的 anomaly 计数，本文更希望保留一个可解释的对照基线。

\section{结果落盘与离线分析}

运行过程中，系统通过 \code{sync} 文件提供实时观察入口，并在收尾阶段生成 \code{data/log} 与 \code{run.json}。其中，\code{data/log} 更适合作为实验章节中的主结果来源，因为它直接体现测试条数、分类计数和代表性 artifact；\code{run.json} 则更适合记录命令行、worker 分片、运行时长、过滤选项和局部统计等元数据。

为了支撑后续人工 triage，当前实现还冻结了 artifact bundle schema 和 second-decoder wrapper 相关约定。这样做的目的不是把所有复杂性都提前自动化，而是保证在发现有争议的样本后，研究者可以追溯到足够完整的上下文，而不是只剩下一个最终计数。

离线分析侧主要由 \code{summarize.py} 与 \code{analysis/} 辅助脚本承担。这些工具负责对最终日志做家族聚合、代表样例展示和 CSV 导出。对一篇强调“可审计”和“可扩展”的论文来说，这部分内容并不是附属脚本，而是实验解释链路的一部分。

\section{MIPS 生成器案例}

MIPS64 是本文当前的首个生成器案例。配置文件 \code{arch-specs/mips64el.json} 仅 36 行，却能够驱动生成器产出 497 行的 \code{src/arch\_mips.c} 基线实现，并自动生成包含 3 条人工检查项的 \code{src/arch\_mips.c.MEMCAGE\_TODO.md}。这说明生成器已经能够回答一个现实问题：当研究者第一次引入新 ISA 时，至少可以快速获得一个可编译、可接入统一编排框架的基线后端。

更重要的是，MIPS 案例把自动化边界写得足够具体。自动化已经能生成基础哨兵、黑名单、\ptracebackend\ 接口骨架和 \capstone\ 配置，但寄存器沙箱语义、BREAK 哨兵行为以及 cache/TLB 敏感指令的黑名单收敛仍必须依赖人工加固。也就是说，本文对“可扩展性”的回答不是“自动替代 ISA 工程”，而是“把高重复劳动压缩到规格输入和基线生成，把剩余人工工作暴露得更清楚”。

\section{当前实现边界}

本文在实现章明确写出三条边界。第一，当前自动分类输出尚未稳定覆盖全部 H/D/P/T/X/UNKNOWN 路径，特别是 illegal-at-test 与 known-disassembly exec-fault 样例仍存在日志空洞。第二，当前所有吞吐结果都来自 Layer A，而非真实硬件。第三，生成器产出的 MIPS 后端虽然已经可运行，但其成熟度还不能与现有 RISC-V、AArch64 后端等量齐观。

把这些边界写在实现章而不是只放在结论里，有助于让后续实验结果与实现状态保持一致，也能降低读者对当前稿“已经完整闭环”的误读。
